The 2030 decennial census will trigger the next congressional redistricting cycle.
Based on the Census Bureau's statutory timeline under P.L.\ 94-171, block-level
redistricting data will be released in April 2031, after which state legislatures,
courts, and commissions must complete redistricting by late 2031. This creates
a 6--12 month operational window that is structurally incompatible with beginning
redistricting infrastructure construction after the census data arrives.

This paper argues that \textsc{bisect} redistricting infrastructure for 2030 must
begin construction no later than 2029, and that waiting until 2031 will produce
one of two outcomes: (1) redistricting plans of degraded quality due to unvalidated
algorithms and incomplete data pipelines, or (2) redistricting plans delayed past
constitutional deadlines, triggering judicial intervention. The argument rests on
four empirically grounded claims.

First, the block-group adjacency graphs required for the 43 states that will require
block-group resolution in 2030 (up from 39 in 2020) take 12--18 months to construct,
validate, and certify. Second, the LODES commuting data and ACS demographic data
must be updated to 2028 vintages to reflect the population distribution at the time
of redistricting. Third, projected 2030 seat counts must be validated against
Census Bureau medium-series projections, with particular attention to the 8 states
gaining seats (TX, FL, AZ, NC, CO, MT, OR, ID) and 10 states losing seats. Fourth,
the DIA statutory seed formula---which determines the initial partition seed from
which \textsc{bisect} begins its recursive bisection---is hard-coded in statute for
pre-registered algorithm versions and cannot be changed after census data release
without triggering a judicial transparency challenge.

The Q-series papers (Q.1--Q.4) provide detailed analyses of each preparation task.
This overview paper provides the timeline, task dependencies, and resource estimates
that frame the series.
